\documentclass[11pt]{article}
\usepackage[a4paper,margin=0.88in]{geometry}
\usepackage{amsmath,amssymb,amsthm,mathtools}
\usepackage{graphicx}
\usepackage{microtype}
\usepackage[hidelinks]{hyperref}

\newtheorem{theorem}{Theorem}
\newtheorem{lemma}[theorem]{Lemma}
\newtheorem{remark}[theorem]{Remark}
\newcommand{\F}{\mathbb F}
\newcommand{\E}{\mathbb E}
\newcommand{\R}{\mathbb R}

\title{Explicit BCH short Khinchin sets\\at every even exponent}
\author{Partial-solution packet for arXiv:1301.2382}
\date{11 August 2026}

\begin{document}
\maketitle

\begin{center}
\fbox{\parbox{0.91\linewidth}{\textbf{Candidate substantial partial result.}
For every even $p\ge2$ there is a deterministic set
$V\subset\{-1,1\}^n$ of optimal cardinality $O_p(n^{p/2})$ whose
$p$th moments of all linear forms agree \emph{exactly} with those on the
full cube.  The same construction gives deterministic polynomial-size
short Khinchin sets for every real $p\ge1$.}}
\end{center}

\section{The problem}

Rudelson's survey asks, for fixed $p\ge1$, for a small deterministic set
$V\subset\{-1,1\}^n$ and constants $\alpha_p,\beta_p>0$ such that
\begin{equation}\label{eq:short}
 \alpha_p\|a\|_2
 \le \left(\frac1{|V|}\sum_{x\in V}|\langle a,x\rangle|^p\right)^{1/p}
 \le \beta_p\|a\|_2
 \qquad(a\in\R^n).
\end{equation}
The survey says that deterministic constructions of reasonably small
cardinality are unknown.  Its probabilistic construction uses
$N\ge n^{p/2}$ for $p>2$, and it notes that this order is optimal.

\begin{figure}[ht]
\centering
\includegraphics[width=0.92\linewidth]{figures/open_problem_crop.png}
\caption{The short Khinchin construction problem on page 23 of the
official arXiv PDF.}
\end{figure}

\section{Statement of the result}

Let $B_s$ denote any valid upper constant in the classical Khinchin
inequality at exponent $s$.

\begin{theorem}\label{thm:main}
For every $n\ge1$ the following deterministic constructions exist.
\begin{enumerate}
 \item If $p=2m$ is even, there is
 $V_{n,m}\subset\{-1,1\}^n$ with
 \[
   |V_{n,m}|\le [2(n+1)]^m
 \]
 such that, for every $a\in\R^n$,
 \begin{equation}\label{eq:exact}
  \frac1{|V_{n,m}|}\sum_{x\in V_{n,m}}|\langle a,x\rangle|^{2m}
  =2^{-n}\sum_{x\in\{-1,1\}^n}|\langle a,x\rangle|^{2m}.
 \end{equation}
 Consequently \eqref{eq:short} holds with the classical constants and
 cardinality $O_p(n^{p/2})$.

 \item If $p\ge2$ is arbitrary and $2m$ is the least even integer at
 least $p$, the same $V_{n,m}$ satisfies
 \[
  \|a\|_2\le
  \left(\frac1{|V_{n,m}|}\sum_{x\in V_{n,m}}
       |\langle a,x\rangle|^p\right)^{1/p}
  \le B_{2m}\|a\|_2.
 \]

 \item If $1\le p\le2$, a deterministic set of cardinality at most
 $[2(n+1)]^2$ satisfies
 \[
  \frac1{\sqrt3}\|a\|_2\le
  \left(\frac1{|V|}\sum_{x\in V}|\langle a,x\rangle|^p\right)^{1/p}
  \le\|a\|_2.
 \]
\end{enumerate}
\end{theorem}

Thus the deterministic cardinality has the optimal $n^{p/2}$ order at
every even $p>2$.  The theorem is partial relative to the sharp general
problem: it gives $O(n^2)$ rather than the probabilistic $O(n)$ scale at
$p=1$, and rounds $p/2$ upward for non-even $p>2$.

\section{The BCH-type sign set}

Fix $m\ge1$.  Let $Q=2^r$ be the least power of two with $Q-1\ge n$,
so $Q<2(n+1)$, and choose distinct nonzero
$t_1,\ldots,t_n\in\F_Q$.  Write
$\operatorname{Tr}:\F_Q\to\F_2$ for the field trace.  For a seed
$u=(u_1,\ldots,u_m)\in\F_Q^m$, set
\begin{equation}\label{eq:signs}
 \varepsilon_j(u)=(-1)^{\operatorname{Tr}
  (\sum_{\ell=1}^m u_\ell t_j^{2\ell-1})},
 \qquad 1\le j\le n.
\end{equation}
Let $V_{n,m}$ be the set of distinct sign vectors
$(\varepsilon_1(u),\ldots,\varepsilon_n(u))$ as $u$ varies.  This is an
explicit primitive-BCH orthogonal array in trace coordinates.  Plainly
\begin{equation}\label{eq:size}
 |V_{n,m}|\le Q^m<[2(n+1)]^m.
\end{equation}

\begin{lemma}\label{lem:ind}
The uniform distribution on $V_{n,m}$ is $2m$-wise independent.
\end{lemma}

\begin{proof}
The bit map underlying \eqref{eq:signs} is linear over $\F_2$, so the
uniform seed pushes forward to the uniform distribution on its image.
It remains to prove that every set of at most $2m$ coordinate
functionals is linearly independent.

Otherwise, for some nonempty $S\subset\{1,\ldots,n\}$ with
$s=|S|\le2m$, summing the corresponding bits would vanish for every
seed.  Nondegeneracy of the trace pairing then gives
\begin{equation}\label{eq:odd}
 \sum_{j\in S}t_j^{2\ell-1}=0,
 \qquad 1\le\ell\le m.
\end{equation}
In characteristic two, an even power sum is the square of the power sum
at half its exponent.  Iterating this fact and using \eqref{eq:odd}
shows
\[
 \sum_{j\in S}t_j^h=0,\qquad 1\le h\le2m.
\]
In particular these identities hold for $1\le h\le s$.  But the square
matrix $(t_j^h)_{1\le h\le s,\,j\in S}$ has determinant
\[
 \left(\prod_{j\in S}t_j\right)
 \prod_{i<j,\ i,j\in S}(t_j-t_i)\ne0.
\]
It cannot annihilate the all-ones vector.  This contradiction proves the
claimed independence.
\end{proof}

\section{Proof of Theorem~\ref{thm:main}}

Let $p=2m$.  Expanding
$(\sum_{j=1}^n a_j\varepsilon_j)^{2m}$ produces monomials involving at
most $2m$ distinct signs.  Lemma~\ref{lem:ind} says their expectations
under the uniform measure on $V_{n,m}$ equal their expectations under
independent Rademacher signs.  Since the exponent is even, this proves
the exact identity \eqref{eq:exact}.  The classical Khinchin inequality
and \eqref{eq:size} prove part 1.

For part 2, let $2m$ be the least even integer at least $p$.  Pairwise
independence gives, for
$S_a(x)=\langle a,x\rangle$,
\[
 \|S_a\|_{L_2(V_{n,m})}=\|a\|_2.
\]
Monotonicity of normalized finite-measure norms and part 1 give
\[
 \|a\|_2=\|S_a\|_2\le\|S_a\|_p
 \le\|S_a\|_{2m}\le B_{2m}\|a\|_2.
\]

Finally use $V_{n,2}$ when $1\le p\le2$.  Its second and fourth moments
match the full cube, so
\[
 \|S_a\|_2=\|a\|_2,
 \qquad
 \|S_a\|_4^4
 =3\|a\|_2^4-2\sum_j a_j^4
 \le3\|a\|_2^4.
\]
Log-convexity between $L_1,L_2,L_4$ yields
\[
 \|S_a\|_1\ge\frac{\|S_a\|_2^3}{\|S_a\|_4^2}
 \ge\frac1{\sqrt3}\|a\|_2.
\]
Now $\|S_a\|_1\le\|S_a\|_p\le\|S_a\|_2$, completing the proof.

\section{Verification, novelty, and scope}

The construction is algebraic and uses no numerical evidence.  The key
compression from $O(n^{2m})$ naive polynomial seeds to $O(n^m)$ is the
odd-power BCH trace map; Lemma~\ref{lem:ind} includes the full
Vandermonde verification.  Passing from seeds to distinct sign vectors is
legitimate because a linear map has equal-size fibers.

The run indexes, parsed corpus, exact phrase, and bounded searches through
11 August 2026 for short Khinchin constructions, BCH codes, orthogonal
arrays, $k$-wise independent Rademachers, and explicit Euclidean
subspaces were checked.  Pass and Spektor \cite{PassSpektor} study
Khintchine inequalities under $k$-wise independence, while coding theory
standardly supplies the required arrays; no source directly presenting
the synthesis above as an answer to Rudelson's construction question was
located.  Novelty confidence is moderate-to-low because both ingredients
are classical, although the stated consequence appears absent from the
bounded search.

Human review should focus on the intended meaning of ``reasonably small''
in the survey.  The even-exponent theorem achieves the source's optimal
order and is a sharp solved subcase.  The all-$p$ extension is polynomial
but is not claimed sharp outside even exponents.

\begin{thebibliography}{2}
\bibitem{Rudelson}
M.~Rudelson,
\newblock \emph{Recent developments in non-asymptotic theory of random
matrices},
\newblock arXiv:1301.2382 (2013).

\bibitem{PassSpektor}
B.~Pass and S.~Spektor,
\newblock \emph{On Khintchine type inequalities for $k$-wise independent
Rademacher random variables},
\newblock Statistics and Probability Letters \textbf{132} (2018), 35--39;
arXiv:1708.08775.
\end{thebibliography}

\end{document}
